% ============================================================
% Modelo_de_Datos.tex — Modelo de datos del banco de pruebas
% TFG ITI-TFG-B26-0017 · Qwen3-8B / vLLM / NVIDIA A100
% Compilación: pdflatex Modelo_de_Datos.tex (dos pasadas)
% ============================================================
\documentclass[11pt,a4paper]{article}

\usepackage[utf8]{inputenc}
\usepackage[T1]{fontenc}
\usepackage{lmodern}
\usepackage[spanish,es-tabla,es-nodecimaldot]{babel}
\usepackage[margin=2.2cm]{geometry}
\usepackage{underscore}
\usepackage{longtable,booktabs,array}
\usepackage{xcolor}
\usepackage{amsmath}
\usepackage{enumitem}
\usepackage{pdflscape}
\usepackage{tikz}
\usetikzlibrary{positioning,arrows.meta,fit,calc}
\usepackage[hidelinks]{hyperref}

\definecolor{uti}{HTML}{1F3A5F}
\definecolor{utiLight}{HTML}{E8EEF5}
\definecolor{pkc}{HTML}{B8860B}
\definecolor{fkc}{HTML}{2E7D32}
\definecolor{vista}{HTML}{6A1B9A}

\setlength{\parindent}{0pt}
\sloppy
\setlength{\parskip}{0.55em}
\renewcommand{\arraystretch}{1.18}

\newcommand{\PK}{\textcolor{pkc}{\textbf{PK}}}
\newcommand{\FK}{\textcolor{fkc}{\textbf{FK}}}
\newcommand{\campo}[1]{\texttt{#1}}

% ---------- Entidad del diagrama ER ----------
% \entidad{id}{posición}{título}{filas}{color}
\newcommand{\entidad}[5][uti]{%
  \node[anchor=north, fill=#1, text=white, font=\bfseries\small,
        minimum width=6.0cm, inner sep=3pt] (#2h) at (#3) {#4};
  \node[anchor=north, draw=#1, fill=white, font=\scriptsize,
        text width=5.8cm, inner sep=3pt] (#2b) at (#2h.south) {#5};
  \node[draw=#1, thick, fit=(#2h)(#2b), inner sep=0pt] (#2) {};
}
\newcommand{\fila}[3]{\makebox[1.05cm][l]{#1}\makebox[2.9cm][l]{\ttfamily #2}{\color{gray}#3}\\}

% ---------- Diccionario de datos ----------
\newenvironment{diccionario}{%
  \small
  \begin{longtable}{@{}>{\raggedright\ttfamily}p{3.4cm}>{\raggedright}p{2.2cm}>{\raggedright}p{1.6cm}p{0.9cm}>{\raggedright\arraybackslash}p{5.8cm}@{}}
  \toprule
  \normalfont\textbf{Campo} & \textbf{Tipo} & \textbf{Unidad} & \textbf{Nulo} & \textbf{Descripción}\\
  \midrule\endhead
  \bottomrule\endfoot
}{\end{longtable}}

\begin{document}

% ==========================================================
% PORTADA
% ==========================================================
\begin{titlepage}
\centering
\vspace*{1.5cm}
{\color{uti}\rule{\linewidth}{1.2pt}}\\[0.8em]
{\LARGE\bfseries\color{uti} Modelo de Datos\\[0.3em]
\large del banco de pruebas experimental}\\[0.6em]
{\color{uti}\rule{\linewidth}{1.2pt}}\\[2.2em]
{\large Modelado experimental de las fases prefill y decode de\\
Qwen3-8B en NVIDIA A100 mediante vLLM}\\[3em]

\begin{tabular}{rl}
\textbf{Código:} & ITI-TFG-B26-0017\\
\textbf{Autor:} & José Andrés Ureta Granda\\
\textbf{Tutor:} & Andrès Xavier Rubio Proaño\\
\textbf{Institución:} & Universidad Tecnológica Indoamérica\\
 & Facultad de Ingenierías — Ingeniería en TI\\
\textbf{Documento:} & Modelo de datos, versión 0.1 (borrador)\\
\textbf{Fecha:} & 8 de octubre de 2026\\
\textbf{Relacionado:} & \texttt{SPEC.md} (secciones 8–12)\\
\end{tabular}

\vfill
\includegraphics[width=4.5cm]{../../Logos/LOGO-UTI-1.png}
\end{titlepage}

\tableofcontents
\clearpage

% ==========================================================
\section{Propósito y alcance}
% ==========================================================
Este documento define el modelo de datos que soporta la campaña experimental descrita en \texttt{SPEC.md}. Cubre la información que se captura, cómo se relaciona, con qué tipos y unidades se almacena y qué reglas de integridad deben cumplirse para que los modelos predictivos de prefill y decode sean trazables y reproducibles.

El modelo responde a cuatro necesidades:
\begin{enumerate}[nosep]
  \item \textbf{Contexto reproducible:} registrar el hardware, el software y la configuración de vLLM exactos de cada medición.
  \item \textbf{Mediciones crudas:} tiempos por petición y por token, y telemetría de la GPU y del servidor.
  \item \textbf{Métricas derivadas:} TTFT, TPOT/ITL, throughput, goodput, MFU y MBU por corrida.
  \item \textbf{Modelos ajustados:} coeficientes, error de validación y qué corridas se usaron para entrenar y para probar.
\end{enumerate}

\textbf{Fuera de alcance:} el texto de los prompts y de las respuestas generadas. Los prompts son sintéticos y se regeneran de forma determinista a partir de la semilla y de la longitud, y la calidad de las respuestas no se evalúa.

% ==========================================================
\section{Convenciones}
% ==========================================================
\begin{itemize}[nosep]
  \item \textbf{Nombres:} \texttt{snake_case}, en español y sin tildes. Las tablas van en singular.
  \item \textbf{Identificadores:} \texttt{UUID} (v4, texto) para las entidades que se generan en ejecución (\campo{corrida}, \campo{peticion}, \campo{ajuste_modelo}). Para las entidades de contexto (\campo{entorno}, \campo{config_servidor}) el identificador es el \texttt{SHA-256} truncado a 16 caracteres de su contenido canónico en JSON: dos configuraciones idénticas producen el mismo identificador.
  \item \textbf{Tiempo:} las marcas de tiempo de alta resolución son \texttt{INT64} en nanosegundos de reloj monotónico (\texttt{time.perf_counter_ns}) y llevan el sufijo \texttt{_ns}. Las métricas derivadas van en milisegundos con el sufijo \texttt{_ms}. Las fechas de pared van en \texttt{TIMESTAMP} UTC.
  \item \textbf{Tokens:} siempre contados con el tokenizer oficial de Qwen3-8B.
  \item \textbf{Enumeraciones:} se almacenan como texto en minúsculas con valores cerrados (indicados entre llaves).
  \item \textbf{Tipos lógicos:} \texttt{TEXT}, \texttt{INT32}, \texttt{INT64}, \texttt{REAL} (float64), \texttt{BOOL}, \texttt{TIMESTAMP}, \texttt{JSON} (texto con JSON válido). Se mapean de forma directa a Parquet/Arrow y a DuckDB.
  \item \textbf{Leyenda:} \PK{} = clave primaria; \FK{} = clave foránea; los campos marcados con † son derivados (calculados en el análisis, no medidos).
\end{itemize}

% ==========================================================
\section{Visión general}
% ==========================================================
El modelo se organiza en cuatro capas:

\begin{center}
\begin{tabular}{@{}>{\raggedright}p{2.4cm}>{\raggedright}p{6.4cm}>{\raggedright\arraybackslash}p{6.4cm}@{}}
\toprule
\textbf{Capa} & \textbf{Entidades} & \textbf{Origen}\\
\midrule
Contexto & \campo{entorno}, \campo{config_servidor}, \campo{experimento} & Registradas al inicio de la corrida\\
Ejecución & \campo{corrida}, \campo{peticion}, \campo{evento_token} & Cliente de carga (\texttt{runner.py})\\
Telemetría & \campo{muestra_gpu}, \campo{muestra_vllm} & \texttt{telemetry.py} (muestreo a 100 ms)\\
Modelado & \campo{ajuste_modelo}, \campo{coeficiente}, \campo{ajuste_corrida}; vista \campo{resumen_corrida} & \texttt{analysis/fit.py}, \texttt{load.py}\\
\bottomrule
\end{tabular}
\end{center}

La entidad central es \campo{corrida}: \emph{un punto de la matriz experimental} (combinación de $L_{in}$, $L_{out}$, $B$, $\lambda$ y configuración) \emph{en una repetición concreta}. Todas las mediciones cuelgan de ella.

% ==========================================================
% DIAGRAMA ER (página apaisada)
% ==========================================================
\begin{landscape}
\section{Diagrama entidad–relación}
\label{sec:er}
\begin{center}
\resizebox{!}{0.86\textheight}{%
\begin{tikzpicture}[
  rel/.style={-{Stealth[length=2.5mm]}, thick, draw=gray!70!black},
  card/.style={font=\scriptsize\bfseries, fill=white, inner sep=1pt},
]
% --- Contexto (fila superior) ---
\entidad{entorno}{0,0}{entorno}{%
  \fila{\PK}{env_id}{TEXT}
  \fila{}{gpu_modelo}{TEXT}
  \fila{}{gpu_memoria_gb}{REAL}
  \fila{}{gpu_formato}{TEXT}
  \fila{}{driver_version}{TEXT}
  \fila{}{cuda_version}{TEXT}
  \fila{}{vllm_version}{TEXT}
  \fila{}{torch_version}{TEXT}
  \fila{}{proveedor}{TEXT}
  \fila{}{lockfile_sha256}{TEXT}
  \fila{}{\dots}{}}
\entidad{experimento}{7,0}{experimento}{%
  \fila{\PK}{exp_id}{TEXT}
  \fila{}{nombre}{TEXT}
  \fila{}{fase_objetivo}{TEXT}
  \fila{}{modo_carga}{TEXT}
  \fila{}{descripcion}{TEXT}}
\entidad{config}{14,0}{config\_servidor}{%
  \fila{\PK}{config_id}{TEXT}
  \fila{}{modelo_id}{TEXT}
  \fila{}{modelo_revision}{TEXT}
  \fila{}{dtype}{TEXT}
  \fila{}{max_model_len}{INT32}
  \fila{}{gpu_memory_util}{REAL}
  \fila{}{chunked_prefill}{BOOL}
  \fila{}{max_num_batched_tok}{INT32}
  \fila{}{max_num_seqs}{INT32}
  \fila{}{prefix_caching}{BOOL}
  \fila{}{\dots}{}}

% --- Corrida (centro) ---
\entidad{corrida}{7,-7}{corrida}{%
  \fila{\PK}{run_id}{UUID}
  \fila{\FK}{exp_id}{TEXT}
  \fila{\FK}{env_id}{TEXT}
  \fila{\FK}{config_id}{TEXT}
  \fila{}{l_in}{INT32}
  \fila{}{l_out}{INT32}
  \fila{}{concurrencia}{INT32}
  \fila{}{tasa_llegada}{REAL}
  \fila{}{repeticion}{INT32}
  \fila{}{inicio}{TIMESTAMP}
  \fila{}{fin}{TIMESTAMP}
  \fila{}{estado}{TEXT}}

% --- Telemetría (izquierda) ---
\entidad{mgpu}{0,-7}{muestra\_gpu}{%
  \fila{\PK\,\FK}{run_id}{UUID}
  \fila{\PK}{t_ns}{INT64}
  \fila{}{util_gpu_pct}{REAL}
  \fila{}{mem_usada_mib}{REAL}
  \fila{}{potencia_w}{REAL}
  \fila{}{reloj_sm_mhz}{INT32}
  \fila{}{temperatura_c}{REAL}}
\entidad{mvllm}{0,-13.6}{muestra\_vllm}{%
  \fila{\PK\,\FK}{run_id}{UUID}
  \fila{\PK}{t_ns}{INT64}
  \fila{}{pet_ejecutando}{INT32}
  \fila{}{pet_esperando}{INT32}
  \fila{}{kv_cache_uso_pct}{REAL}
  \fila{}{preemptions_total}{INT64}}

% --- Ejecución (abajo) ---
\entidad{peticion}{7,-13.6}{peticion}{%
  \fila{\PK}{req_id}{UUID}
  \fila{\FK}{run_id}{UUID}
  \fila{}{es_warmup}{BOOL}
  \fila{}{tokens_prompt}{INT32}
  \fila{}{tokens_salida}{INT32}
  \fila{}{t_envio_ns}{INT64}
  \fila{}{t_primer_token_ns}{INT64}
  \fila{}{t_fin_ns}{INT64}
  \fila{}{ttft_ms / tpot_ms †}{REAL}
  \fila{}{estado}{TEXT}}
\entidad{evtok}{14,-13.6}{evento\_token}{%
  \fila{\PK\,\FK}{req_id}{UUID}
  \fila{\PK}{idx_token}{INT32}
  \fila{}{t_ns}{INT64}
  \fila{}{n_tokens_chunk}{INT32}
  \fila{}{itl_ms †}{REAL}}

% --- Modelado (derecha) ---
\entidad{ajcor}{14,-7}{ajuste\_corrida}{%
  \fila{\PK\,\FK}{fit_id}{UUID}
  \fila{\PK\,\FK}{run_id}{UUID}
  \fila{}{particion}{TEXT}}
\entidad{ajuste}{21,-7}{ajuste\_modelo}{%
  \fila{\PK}{fit_id}{UUID}
  \fila{\FK}{env_id}{TEXT}
  \fila{\FK}{config_id}{TEXT}
  \fila{}{fase}{TEXT}
  \fila{}{formula}{TEXT}
  \fila{}{metodo}{TEXT}
  \fila{}{mape_pct}{REAL}
  \fila{}{r2}{REAL}
  \fila{}{rmse_ms}{REAL}}
\entidad{coef}{21,-13.6}{coeficiente}{%
  \fila{\PK\,\FK}{fit_id}{UUID}
  \fila{\PK}{nombre}{TEXT}
  \fila{}{valor}{REAL}
  \fila{}{ic95_inf / sup}{REAL}
  \fila{}{unidad}{TEXT}}

% --- Vista derivada ---
\entidad[vista]{resumen}{21,0}{resumen\_corrida (vista)}{%
  \fila{\PK\,\FK}{run_id}{UUID}
  \fila{}{ttft_p50/p90/p99}{REAL}
  \fila{}{tpot_p50/p99}{REAL}
  \fila{}{itl_p99}{REAL}
  \fila{}{throughput_tok_s}{REAL}
  \fila{}{goodput_req_s}{REAL}
  \fila{}{mfu_pct / mbu_pct}{REAL}
  \fila{}{energia_j_token}{REAL}}

% --- Relaciones (flecha: de la entidad hija hacia la padre) ---
\draw[rel] (corrida.north) -- node[card, pos=0.12]{N} node[card, pos=0.88]{1} (experimento.south);
\draw[rel] (corrida.north west) -- node[card, pos=0.12]{N} node[card, pos=0.88]{1} (entorno.south);
\draw[rel] (corrida.north east) -- node[card, pos=0.12]{N} node[card, pos=0.88]{1} (config.south);
\draw[rel] (mgpu.east) -- node[card, pos=0.15]{N} node[card, pos=0.85]{1} (corrida.west);
\draw[rel] (mvllm.east) -- node[card, pos=0.12]{N} node[card, pos=0.88]{1} (corrida.south west);
\draw[rel] (peticion.north) -- node[card, pos=0.15]{N} node[card, pos=0.85]{1} (corrida.south);
\draw[rel] (evtok.west) -- node[card, pos=0.15]{N} node[card, pos=0.85]{1} (evtok.west -| peticion.east);
\draw[rel] (ajcor.west) -- node[card, pos=0.15]{N} node[card, pos=0.85]{1} (ajcor.west -| corrida.east);
\draw[rel] (ajcor.east) -- node[card, pos=0.15]{N} node[card, pos=0.85]{1} (ajcor.east -| ajuste.west);
\draw[rel] (coef.north) -- node[card, pos=0.15]{N} node[card, pos=0.85]{1} (ajuste.south);
\draw[rel, dashed] (ajuste.north west) -- node[card, pos=0.12]{N} node[card, pos=0.88]{1} (config.south east);
\draw[rel, dashed] (ajuste.east) -- ++(1.0,0) |- ($(entorno.north)+(0,0.9)$) node[card, pos=0.04]{N} -- node[card, pos=0.7]{1} (entorno.north);
\draw[rel, draw=vista, dotted] (resumen.south) -- ++(0,-0.9) -| node[card, pos=0.05]{1} node[card, pos=0.95]{1} ($(corrida.north)+(1.5,0)$);

% Leyenda
\node[anchor=north west, draw=gray, font=\scriptsize, align=left, inner sep=4pt] at (11.2,-17.6) {%
  \textbf{Leyenda}\\[2pt]
  \PK{} clave primaria \quad \FK{} clave foránea\\
  † campo derivado\\
  \tikz\draw[rel] (0,0)--(0.8,0); \ N:1 obligatoria\\
  \tikz\draw[rel,dashed] (0,0)--(0.8,0); \ N:1 de ámbito del ajuste\\
  \tikz\draw[rel,draw=vista,dotted] (0,0)--(0.8,0); \ vista 1:1 derivada};
\end{tikzpicture}}
\end{center}
\end{landscape}

% ==========================================================
\section{Diccionario de datos}
% ==========================================================

\subsection{\texttt{entorno}}
Hardware y software exactos en los que se ejecutó una corrida. Se registra automáticamente desde \texttt{nvidia-smi}, \texttt{pip freeze} y el sistema operativo.
\begin{diccionario}
env_id & TEXT & — & no & \PK{} Hash SHA-256 (16 car.) del JSON canónico de los demás campos.\\
gpu_modelo & TEXT & — & no & Nombre reportado, p.\,ej. \texttt{NVIDIA A100-SXM4-80GB}.\\
gpu_uuid & TEXT & — & no & UUID físico de la GPU; distingue tarjetas del mismo modelo.\\
gpu_memoria_gb & REAL & GB & no & Memoria total de la GPU.\\
gpu_formato & TEXT & — & no & \{pcie, sxm4\}.\\
ancho_banda_pico_gbs & REAL & GB/s & no & Ancho de banda HBM de la hoja de datos (para MBU).\\
flops_pico_tflops & REAL & TFLOPS & no & Pico BF16 Tensor Core denso (para MFU).\\
driver_version & TEXT & — & no & Versión del driver NVIDIA.\\
cuda_version & TEXT & — & no & Versión de CUDA.\\
cpu_modelo & TEXT & — & sí & Modelo de CPU del host.\\
ram_gb & REAL & GB & sí & RAM del host.\\
so & TEXT & — & no & Sistema operativo y kernel.\\
python_version & TEXT & — & no & Versión de Python.\\
vllm_version & TEXT & — & no & Versión exacta de vLLM.\\
torch_version & TEXT & — & no & Versión de PyTorch.\\
imagen_contenedor & TEXT & — & sí & Imagen Docker con digest, si se usa.\\
lockfile_sha256 & TEXT & — & no & Hash del lockfile de dependencias.\\
proveedor & TEXT & — & no & Laboratorio o nube (p.\,ej. \texttt{gcp}, \texttt{runpod}, \texttt{lab-uti}).\\
registrado_en & TIMESTAMP & UTC & no & Primera vez que se observó este entorno.\\
\end{diccionario}

\subsection{\texttt{experimento}}
Catálogo estático de los experimentos E1–E5 definidos en \texttt{SPEC.md} §8.2.
\begin{diccionario}
exp_id & TEXT & — & no & \PK{} \{e1, e2, e3, e4, e5\}.\\
nombre & TEXT & — & no & Nombre corto (p.\,ej. ``Prefill aislado'').\\
fase_objetivo & TEXT & — & no & \{prefill, decode, e2e, interferencia, capacidad\}.\\
modo_carga & TEXT & — & no & \{offline, online\}. \emph{Offline}: lote fijo mediante la API \texttt{LLM}. \emph{Online}: servidor compatible con OpenAI y llegadas de Poisson.\\
descripcion & TEXT & — & sí & Objetivo y diseño del experimento.\\
\end{diccionario}

\subsection{\texttt{config_servidor}}
Parámetros con los que se lanzó vLLM. Cualquier cambio produce un \texttt{config_id} nuevo.
\begin{diccionario}
config_id & TEXT & — & no & \PK{} Hash SHA-256 (16 car.) del JSON canónico de los parámetros.\\
modelo_id & TEXT & — & no & \texttt{Qwen/Qwen3-8B}.\\
modelo_revision & TEXT & — & no & Hash del commit de Hugging Face de los pesos.\\
dtype & TEXT & — & no & \{bfloat16, float16\}. Base: \texttt{bfloat16}.\\
max_model_len & INT32 & tokens & no & Longitud máxima de contexto.\\
gpu_memory_utilization & REAL & fracción & no & Fracción de memoria reservada por vLLM (0–1).\\
chunked_prefill & BOOL & — & no & Factor experimental CP (on/off).\\
max_num_batched_tokens & INT32 & tokens & sí & Tokens máximos por paso del planificador.\\
max_num_seqs & INT32 & secuencias & sí & Secuencias concurrentes máximas.\\
prefix_caching & BOOL & — & no & Debe ser \texttt{false} en toda la campaña.\\
enforce_eager & BOOL & — & no & \texttt{true} desactiva los CUDA graphs.\\
temperatura & REAL & — & no & Muestreo; base 0 (greedy).\\
ignore_eos & BOOL & — & no & \texttt{true}, para forzar la longitud de salida exacta.\\
semilla & INT32 & — & no & Semilla global.\\
args_cli & JSON & — & no & Línea de comandos completa, para auditoría.\\
\end{diccionario}

\subsection{\texttt{corrida}}
Un punto de la matriz experimental en una repetición. Es la unidad mínima sobre la que se calculan los agregados.
\begin{diccionario}
run_id & UUID & — & no & \PK{} Identificador de la corrida.\\
exp_id & TEXT & — & no & \FK{} $\rightarrow$ \texttt{experimento}.\\
env_id & TEXT & — & no & \FK{} $\rightarrow$ \texttt{entorno}.\\
config_id & TEXT & — & no & \FK{} $\rightarrow$ \texttt{config_servidor}.\\
l_in & INT32 & tokens & no & Longitud de entrada objetivo $L_{in}$.\\
l_out & INT32 & tokens & no & Longitud de salida objetivo $L_{out}$ (1 = prefill puro).\\
concurrencia & INT32 & peticiones & no & Tamaño de lote / concurrencia $B$.\\
tasa_llegada & REAL & req/s & sí & $\lambda$ de Poisson; nulo en modo offline.\\
num_peticiones & INT32 & — & no & Peticiones medidas (sin contar el calentamiento).\\
num_warmup & INT32 & — & no & Peticiones de calentamiento descartadas ($\geq 5$).\\
repeticion & INT32 & — & no & Índice de repetición (1..R, $R \geq 5$).\\
semilla_prompts & INT32 & — & no & Semilla del generador de prompts sintéticos.\\
kv_capacidad_tokens & INT64 & tokens & sí & Capacidad de la caché KV reportada por vLLM al iniciar.\\
inicio & TIMESTAMP & UTC & no & Inicio de la corrida.\\
fin & TIMESTAMP & UTC & sí & Fin; nulo si la corrida se interrumpió.\\
estado & TEXT & — & no & \{ok, oom, timeout, error, abortada\}.\\
notas & TEXT & — & sí & Observaciones manuales.\\
\end{diccionario}
\textbf{Unicidad:} (\campo{exp_id}, \campo{env_id}, \campo{config_id}, \campo{l_in}, \campo{l_out}, \campo{concurrencia}, \campo{tasa_llegada}, \campo{repeticion}).

\subsection{\texttt{peticion}}
Una solicitud de inferencia enviada durante una corrida. Las marcas de tiempo se toman en el cliente.
\begin{diccionario}
req_id & UUID & — & no & \PK{}\\
run_id & UUID & — & no & \FK{} $\rightarrow$ \texttt{corrida}.\\
indice & INT32 & — & no & Orden de envío dentro de la corrida (0..n$-$1).\\
es_warmup & BOOL & — & no & \texttt{true} si es una petición de calentamiento (se excluye del análisis).\\
tokens_prompt & INT32 & tokens & no & Tokens reales del prompt (reportados por vLLM en \texttt{usage}).\\
tokens_salida & INT32 & tokens & sí & Tokens generados; nulo si la petición falló.\\
t_envio_ns & INT64 & ns & no & Envío de la petición.\\
t_primer_token_ns & INT64 & ns & sí & Recepción del primer token.\\
t_fin_ns & INT64 & ns & sí & Recepción del último token.\\
ttft_ms † & REAL & ms & sí & $(t_{primer} - t_{envio}) / 10^6$.\\
tpot_ms † & REAL & ms & sí & $(t_{fin} - t_{primer}) / \bigl(10^6 \cdot (\text{tokens\_salida} - 1)\bigr)$; nulo si $\text{tokens\_salida} \le 1$.\\
e2e_ms † & REAL & ms & sí & $(t_{fin} - t_{envio}) / 10^6$.\\
estado & TEXT & — & no & \{ok, error, abortada, timeout\}.\\
codigo_error & TEXT & — & sí & Mensaje o código HTTP del error.\\
\end{diccionario}

\subsection{\texttt{evento_token}}
Cada fragmento (\emph{chunk}) recibido por streaming. Permite calcular la distribución de ITL y detectar pausas causadas por el prefill de otras peticiones (E4).
\begin{diccionario}
req_id & UUID & — & no & \PK{} \FK{} $\rightarrow$ \texttt{peticion}.\\
idx_token & INT32 & — & no & \PK{} Índice acumulado del primer token del fragmento.\\
t_ns & INT64 & ns & no & Instante de recepción en el cliente.\\
n_tokens_chunk & INT32 & tokens & no & Tokens en el fragmento (puede ser $>1$).\\
itl_ms † & REAL & ms & sí & $(t_{ns} - t_{ns}^{\,anterior}) / (10^6 \cdot \text{n\_tokens\_chunk})$; nulo para el primer fragmento.\\
\end{diccionario}

\subsection{\texttt{muestra_gpu}}
Telemetría de la GPU muestreada cada 100 ms mientras dura la corrida (\texttt{nvidia-smi --query-gpu} o DCGM).
\begin{diccionario}
run_id & UUID & — & no & \PK{} \FK{} $\rightarrow$ \texttt{corrida}.\\
t_ns & INT64 & ns & no & \PK{} Instante de la muestra (mismo reloj que el cliente).\\
util_gpu_pct & REAL & \% & no & Utilización de los SM.\\
util_mem_pct & REAL & \% & no & Utilización del controlador de memoria.\\
mem_usada_mib & REAL & MiB & no & Memoria de GPU en uso.\\
potencia_w & REAL & W & no & Consumo instantáneo.\\
reloj_sm_mhz & INT32 & MHz & no & Reloj de los SM (sirve para detectar \emph{throttling}).\\
reloj_mem_mhz & INT32 & MHz & no & Reloj de la memoria.\\
temperatura_c & REAL & °C & no & Temperatura de la GPU.\\
throttle_razones & TEXT & — & sí & Máscara o texto de las razones de \emph{throttling}.\\
\end{diccionario}

\subsection{\texttt{muestra_vllm}}
Estado interno del servidor, obtenido del endpoint Prometheus \texttt{/metrics} de vLLM (solo en modo online).
\begin{diccionario}
run_id & UUID & — & no & \PK{} \FK{} $\rightarrow$ \texttt{corrida}.\\
t_ns & INT64 & ns & no & \PK{} Instante del scrape.\\
pet_ejecutando & INT32 & — & no & Peticiones en ejecución.\\
pet_esperando & INT32 & — & no & Peticiones en cola.\\
kv_cache_uso_pct & REAL & \% & no & Ocupación de la caché KV en la GPU.\\
preemptions_total & INT64 & — & no & Contador acumulado de preemptions.\\
tokens_prompt_total & INT64 & tokens & no & Contador acumulado de tokens de prompt.\\
tokens_generados_total & INT64 & tokens & no & Contador acumulado de tokens generados.\\
\end{diccionario}

\subsection{\texttt{ajuste_modelo}}
Un modelo predictivo ajustado para una fase, válido para un entorno y una configuración.
\begin{diccionario}
fit_id & UUID & — & no & \PK{}\\
fase & TEXT & — & no & \{prefill, decode, e2e, throughput\}.\\
env_id & TEXT & — & no & \FK{} $\rightarrow$ \texttt{entorno}. El modelo solo es válido en ese hardware.\\
config_id & TEXT & — & no & \FK{} $\rightarrow$ \texttt{config_servidor}.\\
formula & TEXT & — & no & Forma funcional, p.\,ej. \texttt{a0 + a1*B*L + a2*B*L\^{}2}.\\
metodo & TEXT & — & no & \{ols, nnls, wls, por\_tramos\}.\\
variable_respuesta & TEXT & — & no & p.\,ej. \texttt{ttft_ms}, \texttt{t_paso_ms}.\\
semilla_particion & INT32 & — & no & Semilla de la partición entrenamiento/prueba.\\
frac_entrenamiento & REAL & fracción & no & 0,7–0,8.\\
n_entrenamiento & INT32 & corridas & no & Corridas usadas para ajustar.\\
n_prueba & INT32 & corridas & no & Corridas \emph{hold-out}.\\
mape_pct & REAL & \% & no & Error porcentual absoluto medio en prueba.\\
r2 & REAL & — & no & Coeficiente de determinación en prueba.\\
rmse_ms & REAL & ms & no & Raíz del error cuadrático medio en prueba.\\
rango_validez & JSON & — & no & Rangos de $L_{in}$, $B$ y $C$ cubiertos por el entrenamiento.\\
creado_en & TIMESTAMP & UTC & no & Fecha del ajuste.\\
\end{diccionario}

\subsection{\texttt{coeficiente}}
\begin{diccionario}
fit_id & UUID & — & no & \PK{} \FK{} $\rightarrow$ \texttt{ajuste_modelo}.\\
nombre & TEXT & — & no & \PK{} p.\,ej. \texttt{a0}, \texttt{a1}, \texttt{b2}.\\
valor & REAL & var. & no & Valor estimado.\\
error_estandar & REAL & var. & sí & Error estándar de la estimación.\\
ic95_inf & REAL & var. & sí & Límite inferior del IC 95\,\%.\\
ic95_sup & REAL & var. & sí & Límite superior del IC 95\,\%.\\
unidad & TEXT & — & no & p.\,ej. \texttt{ms}, \texttt{ms/token}, \texttt{ms/token\^{}2}.\\
\end{diccionario}

\subsection{\texttt{ajuste_corrida}}
Tabla puente N:M que registra qué corridas se usaron en cada ajuste y en qué partición.
\begin{diccionario}
fit_id & UUID & — & no & \PK{} \FK{} $\rightarrow$ \texttt{ajuste_modelo}.\\
run_id & UUID & — & no & \PK{} \FK{} $\rightarrow$ \texttt{corrida}.\\
particion & TEXT & — & no & \{entrenamiento, prueba\}.\\
\end{diccionario}

\subsection{Vista \texttt{resumen_corrida}}
\textcolor{vista}{\textbf{Vista derivada}} (se materializa en Parquet y se recalcula con \texttt{analysis/load.py}). Contiene una fila por corrida con estado \texttt{ok}, calculada solo con peticiones \texttt{ok} y sin calentamiento.
\begin{diccionario}
run_id & UUID & — & no & \PK{} \FK{} $\rightarrow$ \texttt{corrida}.\\
n_ok & INT32 & — & no & Peticiones válidas.\\
ttft_media / p50 / p90 / p99 & REAL & ms & no & Estadísticos de \texttt{ttft_ms}.\\
tpot_media / p50 / p99 & REAL & ms & sí & Estadísticos de \texttt{tpot_ms}.\\
itl_p50 / p99 & REAL & ms & sí & Percentiles de \texttt{itl_ms} sobre todos los eventos.\\
e2e_p50 / p99 & REAL & ms & no & Percentiles de \texttt{e2e_ms}.\\
throughput_tok_s & REAL & tok/s & no & Tokens de salida por segundo (ec.~\ref{eq:thr}).\\
throughput_req_s & REAL & req/s & no & Peticiones completadas por segundo.\\
goodput_req_s & REAL & req/s & sí & Peticiones/s que cumplen el SLO (ec.~\ref{eq:good}).\\
mfu_pct & REAL & \% & sí & Utilización de FLOPs; solo E1 (ec.~\ref{eq:mfu}).\\
mbu_pct & REAL & \% & sí & Utilización de ancho de banda; solo E2 (ec.~\ref{eq:mbu}).\\
util_gpu_media & REAL & \% & no & Media de \texttt{muestra_gpu.util_gpu_pct}.\\
potencia_media_w & REAL & W & no & Media de \texttt{muestra_gpu.potencia_w}.\\
energia_j_token & REAL & J/tok & no & Energía por token generado (ec.~\ref{eq:energia}).\\
preemptions & INT64 & — & sí & Diferencia del contador durante la corrida.\\
\end{diccionario}

% ==========================================================
\section{Métricas derivadas}
% ==========================================================
Sea $\mathcal{P}$ el conjunto de peticiones válidas de una corrida (estado \texttt{ok} y sin calentamiento). Las constantes del modelo son: $N \approx 8{,}2\times10^{9}$ parámetros; $W \approx 16{,}4$\,GB de pesos en BF16; $n_{capas}=36$; $d=4096$; caché KV de $k_{kv} = 2\cdot36\cdot8\cdot128\cdot2 = 147\,456$\,B por token.

\begin{align}
\text{throughput\_tok\_s} &= \frac{\sum_{p\in\mathcal{P}} \text{tokens\_salida}_p}{\bigl(\max_p t_{fin} - \min_p t_{envio}\bigr)/10^{9}} \label{eq:thr}\\[4pt]
\text{goodput\_req\_s} &= \frac{\bigl|\{p\in\mathcal{P} : \text{ttft}_p \le \text{SLO}_{ttft} \wedge \text{tpot}_p \le \text{SLO}_{tpot}\}\bigr|}{\bigl(\max_p t_{fin} - \min_p t_{envio}\bigr)/10^{9}} \label{eq:good}\\[4pt]
\text{MFU} &= \frac{B\left(2N L_{in} + 2\, n_{capas} L_{in}^{2}\, d\right)}{T_{prefill}\cdot P_{pico}} \times 100 \label{eq:mfu}\\[4pt]
\text{MBU} &= \frac{W + B\cdot \bar{C}\cdot k_{kv}}{T_{paso}\cdot BW_{pico}} \times 100 \label{eq:mbu}\\[4pt]
\text{energia\_j\_token} &= \frac{\sum_{i} P_i\,\Delta t_i}{\sum_{p\in\mathcal{P}} \text{tokens\_salida}_p} \label{eq:energia}
\end{align}

En (\ref{eq:mfu}), $T_{prefill}$ se aproxima con la mediana de TTFT de E1, ya que $L_{out}=1$, y el término de atención causal es una aproximación de primer orden. En (\ref{eq:mbu}), $T_{paso}$ es la mediana de TPOT y $\bar{C}$ es el contexto medio durante la generación ($L_{in} + L_{out}/2$). $P_{pico}$ y $BW_{pico}$ se toman de \texttt{entorno}. Los umbrales SLO se fijan en \texttt{SPEC.md} §16 (pendiente).

% ==========================================================
\section{Reglas de integridad y calidad}
% ==========================================================
\begin{enumerate}[label=\textbf{R\arabic*.}, leftmargin=*]
  \item \textbf{Integridad referencial:} toda FK debe existir en la tabla padre. No se permiten huérfanos en \texttt{peticion}, \texttt{evento_token} ni en las muestras.
  \item \textbf{Orden temporal:} $t_{envio} \le t_{primer\,token} \le t_{fin}$, y los \texttt{t_ns} de \texttt{evento_token} son crecientes por petición.
  \item \textbf{Longitudes exactas:} en peticiones \texttt{ok}, $\text{tokens\_prompt} = l_{in}$ y $\text{tokens\_salida} = l_{out}$ (posible gracias a \texttt{ignore_eos}). Una diferencia invalida la corrida.
  \item \textbf{Consistencia de streaming:} $\sum \text{n\_tokens\_chunk} = \text{tokens\_salida}$ por petición.
  \item \textbf{Calentamiento:} las peticiones con \texttt{es_warmup = true} nunca entran en agregados ni ajustes.
  \item \textbf{Corridas válidas:} solo las corridas con \texttt{estado = ok} y sin preemptions alimentan \texttt{resumen_corrida} y los ajustes (las que tienen preemptions se analizan aparte en E5).
  \item \textbf{Aislamiento:} si durante la corrida \texttt{muestra_gpu} registra memoria de procesos ajenos o \emph{throttling} térmico, la corrida se marca en \texttt{notas} y se repite.
  \item \textbf{Prefix caching:} \texttt{config_servidor.prefix_caching = false} en toda la campaña. Los prompts de una misma corrida deben ser distintos entre sí (semilla + índice).
  \item \textbf{Particiones disjuntas:} dentro de un \texttt{fit_id}, ninguna combinación de factores aparece a la vez en entrenamiento y en prueba (todas sus repeticiones caen en la misma partición).
  \item \textbf{Inmutabilidad:} los datos crudos (\texttt{data/raw}) son de solo escritura una vez. Las correcciones se hacen repitiendo la corrida, nunca editando archivos.
\end{enumerate}

% ==========================================================
\section{Almacenamiento físico}
% ==========================================================
Los datos se persisten en \textbf{Apache Parquet} (columnar, con tipos y compresión \texttt{zstd}) y se consultan con pandas o DuckDB. No se requiere un servidor de base de datos.

\begin{center}
\begin{minipage}{0.9\linewidth}
\small
\begin{verbatim}
data/
  raw/                              # inmutable, una carpeta por corrida
    <exp_id>/<run_id>/
      manifest.json                 # corrida + config_servidor + entorno
      peticiones.parquet
      eventos_token.parquet
      muestras_gpu.parquet
      muestras_vllm.parquet         # solo modo online
      vllm_server.log               # log completo del servidor
  processed/                        # regenerable con analysis/load.py
    entorno.parquet
    experimento.parquet
    config_servidor.parquet
    corrida.parquet
    peticion.parquet                # particionado por exp_id
    evento_token.parquet            # particionado por exp_id
    muestra_gpu.parquet
    muestra_vllm.parquet
    resumen_corrida.parquet         # vista materializada
  modelos/                          # regenerable con analysis/fit.py
    ajuste_modelo.parquet
    coeficiente.parquet
    ajuste_corrida.parquet
\end{verbatim}
\end{minipage}
\end{center}

\textbf{Volumen estimado:} la tabla más grande es \texttt{evento_token}. Con una matriz de unos 1\,500–3\,000 corridas, $\sim$100 peticiones por corrida y hasta 1\,024 tokens por petición, se esperan del orden de $10^{7}$–$10^{8}$ filas, unos pocos GB en Parquet comprimido. Esto es manejable en una sola máquina. Si el volumen fuera excesivo, en E1 y E5 se puede guardar solo el primer y el último evento.

\textbf{Versionado:} \texttt{data/raw} se respalda fuera del repositorio Git (es demasiado grande). En el repositorio se versionan los \texttt{manifest.json}, \texttt{processed/resumen_corrida.parquet} y \texttt{modelos/}, que son pequeños y suficientes para reproducir las figuras de la memoria.

% ==========================================================
\section{Flujo de datos y linaje}
% ==========================================================
\begin{center}
\begin{tikzpicture}[
  box/.style={draw=uti, fill=utiLight, rounded corners=2pt, font=\small, align=center, minimum height=1.0cm, text width=2.5cm},
  store/.style={draw=uti, fill=white, font=\small, align=center, minimum height=1.0cm, text width=2.5cm, double},
  ar/.style={-{Stealth}, thick, draw=uti},
  node distance=0.55cm and 0.55cm]
\node[box] (cfg) {config YAML\\(matriz)};
\node[box, right=of cfg] (run) {runner.py\\server + client\\+ telemetry};
\node[store, right=of run] (raw) {data/raw\\(inmutable)};
\node[box, right=of raw] (load) {load.py\\validación R1–R10};
\node[store, below=of load] (proc) {data/processed\\+ resumen};
\node[box, left=of proc] (fit) {fit.py\\roofline.py};
\node[store, left=of fit] (mod) {data/modelos};
\node[box, left=of mod] (fig) {figures.py\\$\rightarrow$ memoria\\LaTeX};
\draw[ar] (cfg) -- (run);
\draw[ar] (run) -- (raw);
\draw[ar] (raw) -- (load);
\draw[ar] (load) -- (proc);
\draw[ar] (proc) -- (fit);
\draw[ar] (fit) -- (mod);
\draw[ar] (mod) -- (fig);
\draw[ar] (proc.south) -- ++(0,-0.5) -| (fig.south);
\end{tikzpicture}
\end{center}

Cada figura o tabla de la memoria incluye en su campo \emph{Fuente} los \texttt{run_id} o el \texttt{fit_id} de origen (RNF4 del SPEC), con lo que se cierra el linaje desde el dato crudo hasta el documento.

% ==========================================================
\section{Ejemplo de registros}
% ==========================================================
Valores ilustrativos (no son mediciones reales).

\small
\textbf{corrida}\\[2pt]
\begin{tabular}{@{}lllrrrrll@{}}
\toprule
run\_id & exp\_id & config\_id & l\_in & l\_out & B & rep. & estado\\
\midrule
\texttt{3f2a…} & e1 & \texttt{9c01…} & 2048 & 1 & 8 & 1 & ok\\
\texttt{77b0…} & e2 & \texttt{9c01…} & 128 & 1024 & 64 & 3 & ok\\
\bottomrule
\end{tabular}

\vspace{0.8em}
\textbf{peticion}\\[2pt]
\begin{tabular}{@{}lllrrrr@{}}
\toprule
req\_id & run\_id & warmup & tok\_prompt & tok\_salida & ttft\_ms & tpot\_ms\\
\midrule
\texttt{a91e…} & \texttt{3f2a…} & false & 2048 & 1 & 412,7 & —\\
\texttt{c4d2…} & \texttt{77b0…} & false & 128 & 1024 & 61,3 & 14,9\\
\bottomrule
\end{tabular}

\vspace{0.8em}
\textbf{ajuste\_modelo} y \textbf{coeficiente}\\[2pt]
\begin{tabular}{@{}llllrr@{}}
\toprule
fit\_id & fase & fórmula & método & MAPE (\%) & $R^2$\\
\midrule
\texttt{e5f8…} & prefill & \texttt{a0 + a1·B·L + a2·B·L²} & nnls & 6,4 & 0,982\\
\bottomrule
\end{tabular}
\quad
\begin{tabular}{@{}lrl@{}}
\toprule
nombre & valor & unidad\\
\midrule
a0 & 18,2 & ms\\
a1 & 0,051 & ms/token\\
a2 & 2,1e-6 & ms/token²\\
\bottomrule
\end{tabular}
\normalsize

% ==========================================================
\section{Decisiones abiertas}
% ==========================================================
\begin{itemize}[nosep]
  \item Umbrales SLO para el goodput (\texttt{SPEC.md} §16).
  \item Si en E1 y E5 se guardan todos los \texttt{evento_token} o solo el primero y el último (por volumen).
  \item Periodo de muestreo de la telemetría (100 ms propuesto; DCGM permite resoluciones más finas).
  \item Destino del respaldo de \texttt{data/raw} (disco institucional, Google Drive o almacenamiento en la nube).
\end{itemize}

\end{document}
